\section{Transforms, cumulants, Stein's identity, and effective characterizations}\label{sec:transforms}

\subsection{Moment-generating and characteristic functions}

\begin{theorem}[MGF, characteristic function, and Fourier transform]\label{thm:transforms}
For $X\sim N(\mu,\sigma^2)$ and represented real $t$,
\begin{align}
 M_X(t)&=\E e^{tX}=\exp\!\left(\mu t+\frac12\sigma^2t^2\right),\label{eq:mgf}\\
 \chi_X(t)&=\E e^{itX}=\exp\!\left(i\mu t-\frac12\sigma^2t^2\right).
 \label{eq:cf}
\end{align}
With the Fourier convention $\widehat f(\xi)=\int f(x)e^{-2\pi i\xi x}\dd x$,
\begin{equation}\label{eq:fourier}
 \widehat{f_{\mu,\sigma}}(\xi)
 =e^{-2\pi i\mu\xi}e^{-2\pi^2\sigma^2\xi^2}.
\end{equation}
All three identities are uniform on bounded parameter boxes; the MGF is uniform on every bounded $t$ interval.
\end{theorem}

\begin{proof}
For the MGF, complete the square by an exact polynomial identity:
\[
 tx-\frac{(x-\mu)^2}{2\sigma^2}
 =-\frac{(x-\mu-\sigma^2t)^2}{2\sigma^2}
   +\mu t+\frac12\sigma^2t^2.
\]
Factor out the constant and translate the normalized Gaussian integral.  On a finite interval this is a rational identity between exponential evaluators; the shifted tails are controlled uniformly by \eqref{eq:poly-gaussian-tail}.

For the characteristic function one may repeat the argument with complex rational pairs, but a fully real proof is cleaner.  Let
\[
 C(t)=\int\phi(x)\cos(tx)\dd x,\qquad
 S(t)=\int\phi(x)\sin(tx)\dd x.
\]
Oddness gives $S(t)=0$.  Differentiate the finite integral, use $x\phi=-\phi'$, and integrate by parts to obtain $C'(t)=-tC(t)$ with $C(0)=1$.  The ODE uniqueness argument of \cref{prop:gaussian-odes} gives $C(t)=e^{-t^2/2}$.  Translation and scaling give \eqref{eq:cf}; substituting $t=2\pi\xi$ gives \eqref{eq:fourier}.  This is the same quantitative architecture as the Gaussian transform proved in the ProveIt zeta packet.
\end{proof}

\begin{corollary}[Self-reciprocal Gaussian]\label{cor:self-fourier}
Under the convention of \eqref{eq:fourier}, $x\mapsto e^{-\pi x^2}$ is fixed by the Fourier transform.  More generally
\[
 \int_{-\infty}^{\infty}e^{-\pi a x^2}e^{-2\pi ibx}\dd x
 =a^{-1/2}e^{-\pi b^2/a}
\]
for every positive represented $a$.
\end{corollary}

\begin{proof}
Scale \eqref{eq:fourier}, or solve the same transform ODE with the normalization obtained from \eqref{eq:gaussian-normalization}.  The positive inverse square root uses the supplied apartness witness for $a$.
\end{proof}

\subsection{Cumulants and Wick's formula}

\begin{corollary}[Cumulants]\label{cor:cumulants}
The cumulant-generating function is
\begin{equation}\label{eq:cgf}
 K_X(t)=\log M_X(t)=\mu t+\frac12\sigma^2t^2.
\end{equation}
Hence
\[
 \kappa_1=\mu,\qquad \kappa_2=\sigma^2,
 \qquad \kappa_n=0\quad(n\ge3).
\]
This proves the normal distribution's vanishing higher cumulants.  It does not by itself prove the full Marcinkiewicz characterization of all distributions with only finitely many nonzero cumulants.
\end{corollary}

\begin{proof}
On a bounded $t$ interval $M_X(t)$ has an explicit positive lower bound, so the represented logarithm is defined.  Apply it to \eqref{eq:mgf} and compare finite Taylor coefficients.  Derivatives of a quadratic vanish after order two by exact arithmetic.
\end{proof}

\begin{theorem}[Finite Wick--Isserlis formula]\label{thm:wick}
Let $(X_1,\dots,X_d)$ be a centered jointly Gaussian vector with covariance matrix $\Sigma$.  For every finite index list $i_1,\dots,i_n$,
\[
 \E\prod_{r=1}^nX_{i_r}=0\quad\text{if $n$ is odd},
\]
and for $n=2m$,
\begin{equation}\label{eq:wick}
 \E\prod_{r=1}^{2m}X_{i_r}
 =\sum_{\text{pairings }P}\prod_{\{r,s\}\in P}\Sigma_{i_r i_s}.
\end{equation}
The sum is over the finite primitive-recursively enumerable set of pairings.
\end{theorem}

\begin{proof}
The multivariate MGF is $\exp(\tfrac12 t^T\Sigma t)$.  Fix the finite order $n$, replace the exponential by its Taylor polynomial through degree $n$, and extract the coefficient of $t_{i_1}\cdots t_{i_n}$.  Odd total degree is absent.  In degree $2m$, choosing $m$ quadratic factors is exactly choosing a pairing; the factorials cancel the ordering multiplicities.  This is a finite polynomial identity.  Alternatively, repeated multivariate integration by parts gives the same pairing recurrence.
\end{proof}

\subsection{Stein's identity and characterization}

\begin{theorem}[Stein identity]\label{thm:stein}
Let $X\sim N(\mu,\sigma^2)$.  Suppose $g$ has an effective derivative on every bounded interval and a primitive-recursive growth bound making $g(x)f_{\mu,\sigma}(x)$ vanish at both tails.  Then
\begin{equation}\label{eq:stein}
 \E[(X-\mu)g(X)]=\sigma^2\E[g'(X)].
\end{equation}
\end{theorem}

\begin{proof}
The density derivative satisfies
$f'_{\mu,\sigma}=-(x-\mu)f_{\mu,\sigma}/\sigma^2$.  Therefore on $[-T,T]$,
\[
 \int (x-\mu)g(x)f(x)\dd x
 =-\sigma^2\int g(x)f'(x)\dd x
 =\sigma^2\int g'(x)f(x)\dd x
   -\sigma^2[g(x)f(x)]_{-T}^T.
\]
The growth/tail contract makes the boundary smaller than any requested dyadic error.  The remaining two integrals have the ordinary finite quadrature contracts.
\end{proof}

\begin{theorem}[Effective Stein characterization]\label{thm:stein-characterization}
Let $Y$ be an effectively presented probability law with finite second moment.  Assume that for every represented trigonometric test $g_t(x)=e^{itx}$ on every bounded $t$ interval,
\begin{equation}\label{eq:stein-assumption}
 \E[(Y-\mu)g_t(Y)]=\sigma^2\E[g_t'(Y)]
\end{equation}
holds with a uniform primitive-recursive modulus.  Then the characteristic function of $Y$ equals \eqref{eq:cf}; consequently all finite interval probabilities agree effectively with $N(\mu,\sigma^2)$.
\end{theorem}

\begin{proof}
Let $\chi(t)=\E e^{itY}$.  Equation \eqref{eq:stein-assumption} becomes
\[
 \frac1i\chi'(t)-\mu\chi(t)=i\sigma^2t\chi(t),
\]
after differentiating under the effective second-moment bound.  Thus
$\chi'(t)=(i\mu-\sigma^2t)\chi(t)$ and $\chi(0)=1$.  Multiplying by
$e^{-i\mu t+\sigma^2t^2/2}$ gives a function with zero derivative, hence the normal characteristic function by finite ODE uniqueness.

To pass from characteristic functions to interval probabilities constructively, convolve both laws with the finite triangular smoothing kernel used in the central-limit packet.  Fourier inversion for that compactly supported kernel is a finite integral identity.  Equality of the characteristic functions gives equality of the smoothed CDFs, and the supplied tightness plus anti-concentration modulus lets the smoothing width tend effectively to zero at rational interval endpoints.
\end{proof}

\subsection{What is and is not characterized}

\begin{proposition}[Quadratic log-MGF criterion]\label{prop:quadratic-mgf}
If an effective law has an MGF on a rational neighborhood of zero and its represented logarithm is exactly
$\mu t+\sigma^2t^2/2$ there, then the law is $N(\mu,\sigma^2)$.
\end{proposition}

\begin{proof}
Exponentiate to recover \eqref{eq:mgf}, substitute imaginary arguments through the common finite power series, and apply the characteristic-function uniqueness step above.  Equivalently, equality of all derivatives at zero plus an explicit analytic remainder gives equality of MGFs on a neighborhood and then of characteristic functions.
\end{proof}

\begin{remark}[Deep classical characterizations]
The stronger Marcinkiewicz theorem (a polynomial log characteristic function has degree at most two), Cramer's decomposition theorem, the Kac--Bernstein theorem, and Geary's characterization by independence of sample mean and variance require substantial additional harmonic or functional-equation arguments.  They are listed as D in the ledger.  This package proves every forward Gaussian statement used by those theorems, and the analytic quadratic-MGF criterion, but does not disguise those deeper converses as elementary consequences of cumulant algebra.
\end{remark}
